% =====================================================================
\chapter{Exponenciales y logaritmos}
\label{tema:5}
% =====================================================================

El logaritmo convierte productos en sumas y potencias en productos. Esa propiedad lo hace omnipresente en el curso: la log-verosimilitud, la entropía, la perplejidad y la escala de las leyes de potencia son logaritmos. Este tema repasa las propiedades de exponenciales y logaritmos siguiendo el capítulo~6 de OpenStax~\cite{openstax}, introduce la función log-suma-exp siguiendo a Blanchard, Higham y Higham~\cite{blanchard} y fija las unidades de información (bits y nats) como en el capítulo~2 de Cover y Thomas~\cite{ct}. La construcción rigurosa de $e^x$ y $\ln x$ requiere el cálculo integral~\cite[cap.~18]{spivak}; aquí se admiten sus propiedades básicas.

\section{Potencias de exponente real y la función exponencial}

Para $a>0$, la potencia $a^r$ está definida para $r$ racional (\cref{tema:1}) y se extiende a todo $x\in\R$ de forma que la función $x\mapsto a^x$ es continua. Se admiten las siguientes propiedades, que extienden la \cref{prop:exponentes}.

\begin{proposicion}[Función exponencial]\label{prop:exp}
Sea $a>0$. Para todos $x,y\in\R$: $a^x>0$; $a^{x+y}=a^xa^y$; $(a^x)^y=a^{xy}$; $a^{-x}=1/a^x$. Si $a>1$, la función $x\mapsto a^x$ es estrictamente creciente; si $0<a<1$, es estrictamente decreciente. Si $a\neq1$, toma cada valor positivo exactamente una vez.
\end{proposicion}

El número $e=\lim_{n\to\infty}(1+\frac1n)^n\approx2{,}71828$ es la base natural; más en general, $e^x=\lim_{n\to\infty}(1+\frac xn)^n$ para todo $x$~\cite[cap.~18]{spivak}.

\begin{proposicion}\label{prop:ex-cota}
Para todo $x\in\R$, $e^x\ge1+x$.
\end{proposicion}

\begin{proof}
Si $x\le-1$, el miembro derecho es $\le0<e^x$. Si $x>-1$, para todo $n\in\N$ se tiene $x/n>-1$, y la desigualdad de Bernoulli (\cref{teo:bernoulli}) da $(1+\frac xn)^n\ge1+x$. Pasando al límite, que conserva las desigualdades no estrictas, $e^x\ge1+x$.
\end{proof}

\section{Logaritmos}

\begin{definicion}[Logaritmo]\label{def:log}
Sean $a>0$, $a\neq1$ y $x>0$. El \emph{logaritmo en base $a$} de $x$ es el único $y\in\R$ tal que $a^y=x$; se escribe $y=\log_ax$. Se usan $\ln=\log_e$ (logaritmo natural), $\log_2$ (binario) y $\log_{10}$ (decimal).
\end{definicion}

Por definición, $a^{\log_ax}=x$ para $x>0$ y $\log_a(a^y)=y$ para todo $y$. El logaritmo solo está definido para argumentos positivos.

\begin{teorema}[Propiedades de los logaritmos]\label{teo:log}
Sean $a,b>0$ distintos de~1, $x,y>0$ y $r\in\R$. Entonces:
\begin{enumerate}
  \item $\log_a(xy)=\log_ax+\log_ay$ y $\log_a(x/y)=\log_ax-\log_ay$;
  \item $\log_a(x^r)=r\log_ax$;
  \item cambio de base: $\log_ax=\dfrac{\log_bx}{\log_ba}$; en particular, $\log_ab\cdot\log_ba=1$;
  \item si $a>1$, $x<y$ si y solo si $\log_ax<\log_ay$.
\end{enumerate}
\end{teorema}

\begin{proof}
Sean $u=\log_ax$ y $v=\log_ay$, es decir, $a^u=x$ y $a^v=y$. (i) $a^{u+v}=a^ua^v=xy$, luego $\log_a(xy)=u+v$; el cociente es análogo con $a^{u-v}=x/y$. (ii) $a^{ru}=(a^u)^r=x^r$, luego $\log_a(x^r)=ru$. (iii) Tomando $\log_b$ en $a^u=x$ y usando (ii): $u\log_ba=\log_bx$; como $a\neq1$, $\log_ba\neq0$ y se despeja $u$. Con $x=b$ se obtiene $\log_ab=1/\log_ba$. (iv) Por la \cref{prop:exp}, $t\mapsto a^t$ es estrictamente creciente si $a>1$, así que $u<v$ si y solo si $a^u<a^v$.
\end{proof}

\begin{corolario}\label{cor:lnx}
Para todo $x>0$, $\ln x\le x-1$, con igualdad si y solo si $x=1$. En consecuencia, $\ln x\ge1-\frac1x$.
\end{corolario}

\begin{proof}
Por la \cref{prop:ex-cota} con $\ln x$ en lugar de $x$: $x=e^{\ln x}\ge1+\ln x$. Aplicándolo a $1/x$: $-\ln x=\ln\frac1x\le\frac1x-1$. (La igualdad solo en $x=1$ se sigue de que $e^t>1+t$ para $t\neq0$, que se demuestra en el módulo de cálculo.)
\end{proof}

\section{Ecuaciones exponenciales y logarítmicas}

Para resolver una ecuación exponencial se toman logaritmos o se cambia la incógnita ($t=a^x$); en las logarítmicas se usan las propiedades del \cref{teo:log} para llegar a $\log_aA=\log_aB$, que equivale a $A=B$, y se descartan las soluciones que hacen negativo o nulo algún argumento.

\begin{ejemplo}\label{ej:ecuacion-exp}
$9^x-4\cdot3^x+3=0$. Con $t=3^x>0$, queda $t^2-4t+3=0$, de raíces $t=1$ y $t=3$, que dan $x=0$ y $x=1$.
\end{ejemplo}

\begin{ejemplo}\label{ej:ecuacion-log}
$\log_3(x+1)+\log_3(x-1)=1$. Exige $x>1$. Por el \cref{teo:log}, $\log_3(x^2-1)=1$, es decir, $x^2-1=3$ y $x=\pm2$. Solo $x=2$ cumple $x>1$.
\end{ejemplo}

\section{La función log-suma-exp}

\begin{definicion}[Log-suma-exp]\label{def:lse}
Para $z=(z_1,\dots,z_n)\in\R^n$, $\displaystyle\LSE(z)=\ln\sum_{i=1}^ne^{z_i}$. La función \emph{softmax} es $\softmax(z)_i=e^{z_i}/\sum_je^{z_j}=e^{z_i-\LSE(z)}$.
\end{definicion}

\begin{proposicion}\label{prop:lse}
Para todo $z\in\R^n$ y $c\in\R$, con $M=\max_iz_i$:
\begin{enumerate}
  \item invariancia por traslación: $\LSE(z_1+c,\dots,z_n+c)=\LSE(z)+c$; en consecuencia, la softmax no cambia al sumar $c$ a todas las componentes;
  \item cotas: $M\le\LSE(z)\le M+\ln n$.
\end{enumerate}
\end{proposicion}

\begin{proof}
(i) $\ln\sum_ie^{z_i+c}=\ln\bigl(e^c\sum_ie^{z_i}\bigr)=c+\LSE(z)$ por el \cref{teo:log}(i). Para la softmax, el factor $e^c$ aparece en el numerador y en el denominador. (ii) La suma $\sum_ie^{z_i}$ contiene el término $e^M$ y todos sus términos son positivos y menores o iguales que $e^M$, así que $e^M\le\sum_ie^{z_i}\le ne^M$. Tomando $\ln$, que es creciente, se obtienen las cotas.
\end{proof}

La \cref{prop:lse}(i) con $c=-M$ da el \emph{cálculo estable}: $\LSE(z)=M+\ln\sum_ie^{z_i-M}$. Todas las exponenciales son ahora $\le1$, de modo que no hay desbordamiento, y al menos una vale exactamente~1, de modo que el logaritmo no recibe un cero.

\begin{ejemplo}\label{ej:lse-estable}
Con $z=(1000,1000)$, el cálculo directo de $e^{1000}$ desborda en doble precisión (el mayor número representable es del orden de $1{,}8\cdot10^{308}$, mientras que $e^{1000}\approx10^{434}$). El cálculo estable da $\LSE(z)=1000+\ln(1+1)=1000+\ln2\approx1000{,}693$, y la softmax es $(\frac12,\frac12)$.
\end{ejemplo}

\section{Escalas logarítmicas, bits y nats}

Una cantidad que varía en muchos órdenes de magnitud se representa mejor por su logaritmo. Si $y=Cx^{k}$, entonces $\log y=\log C+k\log x$: en ejes doblemente logarítmicos una \emph{ley de potencia} es una recta de pendiente $k$, que es como se presentan las leyes de escala del módulo~16. Si $y=Ca^{x}$, entonces $\log y$ es lineal en $x$ (escala semilogarítmica).

\begin{definicion}[Unidades de información]
La \emph{información} de un suceso de probabilidad $p\in(0,1]$ es $-\log_2p$ \emph{bits} o $-\ln p$ \emph{nats}~\cite[cap.~2]{ct}. Por el cambio de base, $1\ \text{nat}=\frac1{\ln2}\ \text{bits}\approx1{,}4427$ bits.
\end{definicion}

\begin{ejemplo}[Productos de probabilidades]\label{ej:underflow}
Un texto de 500 tokens en el que el modelo asigna probabilidad $10^{-3}$ a cada token tiene probabilidad $10^{-1500}$, muy por debajo del menor número positivo representable en doble precisión (del orden de $10^{-308}$, o $5\cdot10^{-324}$ con números subnormales): el producto se redondea a~0. Su logaritmo, en cambio, es $500\ln10^{-3}\approx-3453{,}9$ nats, perfectamente representable. Por eso los modelos trabajan con log-probabilidades y la regla de la cadena se aplica como suma de logaritmos.
\end{ejemplo}

\begin{observacion}[Errores frecuentes]
\begin{enumerate}[label=(\alph*)]
\item $\log(x+y)\neq\log x+\log y$; la propiedad es $\log(xy)=\log x+\log y$.
\item $(\ln x)^2\neq\ln(x^2)=2\ln x$, y $\ln x/\ln y\neq\ln(x/y)$.
\item $a^{x+y}=a^xa^y$, no $a^x+a^y$.
\item Al resolver una ecuación logarítmica hay que comprobar que los argumentos son positivos (\cref{ej:ecuacion-log}).
\item Confundir $\log$ (que en muchos libros de aprendizaje automático y en Python significa $\ln$) con $\log_{10}$.
\end{enumerate}
\end{observacion}

\begin{observacion}[Conexión con el curso]
La log-verosimilitud, la entropía cruzada y la perplejidad del módulo~1 son sumas de logaritmos; el cambio de base convierte nats en bits (módulo~6). La función log-suma-exp es el denominador de la softmax en logaritmos: aparece en la entropía cruzada (módulo~8) y en la atención (módulo~13), y las cotas de la \cref{prop:lse} explican por qué la softmax con temperatura tiende al máximo cuando la temperatura tiende a cero (\cref{eje:5-temperatura}). La desigualdad $\ln x\le x-1$ es el paso clave de la desigualdad de Gibbs (módulo~6).
\end{observacion}

\begin{lecturas}
\item OpenStax \emph{College Algebra}~\cite{openstax}, §6.1 «Exponential Functions», §6.3 «Logarithmic Functions», §6.5 «Logarithmic Properties», §6.6 «Exponential and Logarithmic Equations» y §6.7 «Exponential and Logarithmic Models».
\item Spivak~\cite{spivak}, cap.~18 «The Logarithm and Exponential Functions», para la construcción rigurosa.
\item Blanchard, Higham y Higham~\cite{blanchard}, secciones introductorias, sobre el cálculo estable de log-suma-exp y softmax.
\item Cover y Thomas~\cite{ct}, cap.~2, para bits, nats y la desigualdad de la información.
\end{lecturas}

\section*{Ejercicios del tema 5}
\addcontentsline{toc}{section}{Ejercicios del tema 5}

\begin{ejercicio}{1}{eje:5-calcular}
Calcule sin calculadora: $\log_264$, $\log_3\frac1{81}$, $\log_{10}0{,}001$, $\ln e^7$, $\log_48$ y $10^{\log_{10}5}$.
\fuente{adaptado de OpenStax \emph{College Algebra}~\cite{openstax}, §6.3.}
\end{ejercicio}

\begin{ejercicio}{1}{eje:5-vf}
Razone si son verdaderas o falsas para $a,b,x>0$: (a) $\ln(a+b)=\ln a+\ln b$; (b) $\log_2(8^x)=3x$; (c) $\ln\frac1x=-\ln x$; (d) $(\ln x)^2=2\ln x$.
\end{ejercicio}

\begin{ejercicio}{1}{eje:5-unidades}
(a) ¿Cuántos bits son 8 nats? (b) ¿Cuánta información, en bits y en nats, tiene un suceso de probabilidad $1/1024$? (c) ¿Cuántos bits hacen falta como mínimo para numerar los tokens de un vocabulario de $50\,257$ tokens?
\end{ejercicio}

\begin{ejercicio}{2}{eje:5-simplificar}
Simplifique: (a) $\log_212-\log_23+\log_2\frac18$; (b) $\frac12\ln16+\ln\frac e4$; (c) $\log_23\cdot\log_38$; (d) $\log_530$ expresado con logaritmos naturales, y su valor aproximado.
\fuente{adaptado de OpenStax \emph{College Algebra}~\cite{openstax}, §6.5.}
\end{ejercicio}

\begin{ejercicio}{2}{eje:5-exponenciales}
Resuelva: (a) $2^{x+1}=3^x$; (b) $4^x-3\cdot2^x-4=0$; (c) $e^{2x}-5e^x+6=0$.
\fuente{adaptado de OpenStax \emph{College Algebra}~\cite{openstax}, §6.6.}
\end{ejercicio}

\begin{ejercicio}{2}{eje:5-logaritmicas}
Resuelva: (a) $\log_2x+\log_2(x-2)=3$; (b) $\ln(x+1)-\ln(x-1)=\ln3$.
\fuente{adaptado de OpenStax \emph{College Algebra}~\cite{openstax}, §6.6.}
\end{ejercicio}

\begin{ejercicio}{2}{eje:5-inecuaciones}
Resuelva: (a) $0{,}5^x>8$; (b) $\log_3(2x-1)\le2$.
\end{ejercicio}

\begin{ejercicio}{2}{eje:5-lse}
(a) Calcule $\LSE(z)$ y $\softmax(z)$ para $z=(2,1,0,-1)$ y compruebe las cotas de la \cref{prop:lse}. (b) Calcule de forma estable $\LSE(1000,999,998)$.
\end{ejercicio}

\begin{ejercicio}{2}{eje:5-perplejidad}
Un modelo asigna a un texto de $N$ tokens una log-verosimilitud media de $-2{,}3$ nats por token. (a) Exprese esa cantidad en bits por token. (b) Calcule la perplejidad $e^{2{,}3}$ y su equivalente $2^{H}$ con $H$ en bits. (c) ¿Qué probabilidad constante por token daría la misma perplejidad?
\end{ejercicio}

\begin{ejercicio}{2}{eje:5-ley-potencia}
Una pérdida sigue la ley de potencia $L(N)=(N_c/N)^{\alpha}$ con $\alpha=0{,}076$. (a) ¿En qué factor se multiplica $L$ al multiplicar $N$ por~10? (b) ¿Por qué factor hay que multiplicar $N$ para reducir $L$ a la mitad? (c) ¿Qué pendiente tiene la recta $\log L$ frente a $\log N$?
\fuente{adaptado de OpenStax \emph{College Algebra}~\cite{openstax}, §6.7 (modelos exponenciales y logarítmicos).}
\end{ejercicio}

\begin{ejercicio}{3}{eje:5-log23}
Demuestre que $\log_23$ es irracional.
\end{ejercicio}

\begin{ejercicio}{3}{eje:5-temperatura}
Sea $z\in\R^n$ con máximo $M$, y $T>0$. Demuestre que $M\le T\,\LSE(z/T)\le M+T\ln n$ y deduzca que $T\,\LSE(z/T)\to M$ cuando $T\to0^{+}$.
\end{ejercicio}

\begin{ejercicio}{3}{eje:5-lse-media}
Demuestre, usando la desigualdad de las medias, que $\LSE(z)\ge\frac1n\sum_{i=1}^nz_i+\ln n$ para todo $z\in\R^n$. ¿Cuándo se da la igualdad?
\end{ejercicio}

\begin{ejercicio}{3}{eje:5-gibbs}
Sean $p=(p_1,\dots,p_n)$ y $q=(q_1,\dots,q_n)$ distribuciones de probabilidad con todos los $p_i,q_i>0$. Demuestre, usando el \cref{cor:lnx}, que $\sum_ip_i\ln\frac{q_i}{p_i}\le0$, es decir, $-\sum_ip_i\ln q_i\ge-\sum_ip_i\ln p_i$.
\fuente{adaptado de Cover y Thomas~\cite{ct}, cap.~2 (desigualdad de la información).}
\end{ejercicio}

\begin{ejercicio}{3}{eje:5-e-creciente}
Aplicando la desigualdad de las medias a los $n+1$ números $1+\frac1n,\dots,1+\frac1n,\,1$, demuestre que $\bigl(1+\frac1n\bigr)^n<\bigl(1+\frac1{n+1}\bigr)^{n+1}$ para todo $n\in\N$.
\end{ejercicio}
